\section{Introduction}

\subsection{A degeneracy is not self-describing}

When two modes of a chemical system approach the same frequency and linewidth, a spectrum alone cannot say what kind of degeneracy it is. Two possibilities look identical in the eigenvalues but differ completely in the physics. In a semisimple (diagonalizable) degeneracy the two modes share an eigenvalue but keep independent eigenvectors: perturb the system and they separate smoothly. In an exceptional point the eigenvectors themselves coalesce; the system becomes defective, and it responds to perturbation with the square-root sensitivity that makes exceptional points valuable for sensing and dangerous to misidentify. The two are distinguished not by how close the eigenvalues are but by whether the eigenvectors have collapsed, a property a frequency-and-linewidth fit does not report.

A second trap is layered on the first. A degeneracy read from a reduced description, an open-system or effective equation where the environment has been integrated out, is not the same object as one in the full generator. A repeated root in a Lindblad-reduced equation can look like an exceptional point while corresponding to nothing defective in the underlying closed system. Certifying a physical exceptional point therefore requires knowing the provenance of the equation it came from, and many reported degeneracies do not carry that provenance.

\subsection{The machinery first}

The primary contribution is a classification engine that closes both traps, together with a scope fence that keeps it honest. ChemSA distinguishes a defective degeneracy from a semisimple one by testing eigenvector deficiency directly, returning different verdicts for two systems with the same repeated eigenvalue (Section~2.3); it gates the exceptional-point verdict on declared provenance, so defectiveness found in a reduced equation is reported as an exceptional point of that supplied representation with the parent status left explicitly unresolved; and it is bounded by a constructive counterexample, two systems with identical local architecture and activated rates differing by a factor of about 148 under a common unspecified prefactor (Section~3). The current ChemSA interface contains no reaction-rate estimator. The constructive counterexample of Section~3 is the mathematical statement of why local spectral architecture alone cannot supply one; the regression test is a software guard that keeps the implementation from drifting back across that boundary. The proof bounds the claim, the test bounds the code. These three core scientific safeguards, together with three implementation validity gates (scalar-\(\chi\) applicability, source gating of numerical inputs, and finite-precision equality caution; Section~2.3) are the defensible core: for a single published damped-oscillator fit the damping ratio is already known, so a catalogue of such ratios adds little; what the engine adds is the exceptional-point-versus-semisimple discrimination, the provenance gate, and the scope fence, none of which a standard fit provides.

\subsection{The coordinate, and the boundary of where it is a damping ratio}

To place a mode's stability on an axis the engine uses the damping ratio $\chi$ = \(\gamma/2\omega_0\), read from the poles of the linearized generator. What $\chi$ is must be stated carefully, because it is easy to overreach.

The mechanical damping ratio is defined for a second-order modal representation: a scalar \(M,C,K\) mode, or a system whose damping is proportional so that it decouples into independent modes (the Caughey--O'Kelly condition \citep{caughey}). That condition is TESTED by the engine, not accepted on declaration: \texttt{scalar\_reduction\_valid} computes the relative mass-metric commutator of \(M^{-1}C\) and \(M^{-1}K\) and withholds the mechanical \(\chi\), with a stated reason, whenever it is nonzero. Where the test passes, the modal basis is fixed by diagonalizing the damping within each degenerate stiffness block, since a repeated stiffness eigenvalue leaves the basis undetermined and only one rotation inside that subspace decouples the damping as well. In that representation \(\chi = \gamma/2\omega_0\) is a genuine mechanical damping ratio: $\chi$ \(<1\) rings, $\chi$ \(=1\) is critical, $\chi$ \(>1\) is two real decay timescales. This is the regime in which $\chi$ is meaningful as a continuous mechanical coordinate.

For a general first-order generator (an arbitrary kinetic Jacobian, a non-modal coupled system) no mechanical $\chi$ is claimed. What the engine reports there is a spectral pole-angle index \(\rho = -\mathrm{Re}(\lambda)/|\lambda|\) together with a categorical temporal class (oscillatory, repeated-root, or monotone). The distinction matters and a careful reader will insist on it: \(\rho\) is the cosine of the pole's angle, and nothing more. It EQUALS a mechanical damping ratio only when the pole belongs to a scalar or classically damped second-order factor that has been separately licensed; for a general first-order generator there is no mechanical \(\chi\) for it to equal, and reading \(\rho\) as one is the precise conflation this work exists to prevent. For a stable real pole \(\rho = 1\) identically, for every stable monotone system, no matter how far from any critical point, so \(\rho\) cannot by itself distinguish a repeated-root mode from ordinary stable real relaxation; only the second-order modal picture recovers that distinction. Because that value carries no information, ChemSA does not compute or report it: \(\rho\) is returned only for a complex pole, and is omitted (not set to 1) for a real one, so its presence in a record already signals which case applies. Accordingly, ChemSA assigns a true $\chi$ only when a second-order or decoupled modal representation is established, and otherwise reports pole geometry and temporal class without forcing a mechanical $\chi$.

What linearization does give, stated honestly, is a common spectral footing: disparate systems reduce to poles, and the engine reads those poles the same way. That is what lets one demonstration span phases, a property of linearization, not a claim that one number governs disparate physics, and not a claim that arbitrary first-order networks occupy the same continuous under-, critical-, over-damped axis.

\subsection{What this paper does not claim}

This work makes no claim about reaction rate, yield, commitment probability, or selectivity: the current interface exposes no rate estimator, and the scope proof of Section~3 shows why local classification cannot determine rate (identical local architecture, rates differing by 148). It does not claim a real system that rests at $\chi$ \(=1\): the critical point is a boundary systems cross rather than a state a real system occupies. The linewidth hierarchy of Part~III does not claim a coupled-mode mechanical catalogue: its mechanical-placement claim is confined to a single reconstructed trajectory in which an independent lifetime measurement is available, and a chemical system whose coupled generator sits on an exceptional point remains open work, awaiting two linewidths that only a dedicated experiment can supply.

\subsection{Relation to prior work}

The scalar damping coordinate used here is not introduced by this paper. Earlier work established \(\chi = \gamma/2\omega_0\) as the structural boundary of a dominant second-order response, showed that \(\chi = 1\) marks a boundary that systems cross rather than a state they occupy or optimize toward, restricted exceptional-point terminology to defective realizations, and noted that a strongly multimode system cannot be compressed into a single scalar \(\chi\) \citep{stabilityarch}. None of those results is re-claimed here. What the present work adds is that each of them becomes a test that can refuse: the restriction to a scalar response becomes a Caughey--O'Kelly check that withholds \(\chi\) with a stated reason, the restriction of exceptional-point terminology to defective realizations becomes an algebraic-versus-geometric multiplicity test, and the caution that a reduced description does not establish its parent becomes the provenance gate. Section~3 adds one result in the other direction: the boundary is specifically a property of a stable well, and no counterpart exists at an inverted mode.

What the present work adds is the machinery that those statements require once the system is not scalar. In a scalar second-order mode, \(\chi = 1\) produces a defective companion block, so the critical point and the exceptional point coincide and neither has to be tested for separately. In a coupled spectrum that identity fails: an eigenvalue coincidence may be defective or semisimple, and the two are indistinguishable in the eigenvalues alone, so defectiveness becomes something to test rather than something to infer. This paper therefore generalizes the treatment from a dominant scalar mode to coupled quadratic pencils and general first-order generators, and supplies what that generalization needs: an executable algebraic-versus-geometric multiplicity test, a refusal for neighborhoods that finite precision cannot resolve, enforcement of the representation provenance in which a degeneracy was found, and a biorthogonal response layer with channel-declared residues and sensitivities. It further supplies the spectroscopy-side question that the scalar boundary leaves open, namely when a measured linewidth delivers the amplitude-damping coefficient that a mechanical \(\chi\) requires at all.

The relationship is therefore one of operationalization rather than rediscovery. The earlier work stated the boundaries; this work makes them testable, and makes the cases outside them refusable.

\noindent\textbf{Reader's guide.} A reader interested in the classification engine should read Sections~2.2 to 2.5 (vocabulary, safeguards, response layer) and the synthetic model-class cross-check in Section~2.6. A reader interested in reading a linewidth as a damping ratio should read Sections~5 and 6. The claim-status table in the Supplementary Material records what is established, what is demonstrated, and what is explicitly not claimed.
